\let\R\relax
\newcommand{\R}{\mathbb{R}}
\let\C\relax
\newcommand{\C}{\mathbb{C}}
\let\N\relax
\newcommand{\N}{\mathbb{N}}
\let\Z\relax
\newcommand{\Z}{\mathbb{Z}}
\let\Q\relax
\newcommand{\Q}{\mathbb{Q}}
\let\SO\relax
\newcommand{\SO}{\operatorname{SO}}
\let\SL\relax
\newcommand{\SL}{\operatorname{SL}}
\let\GL\relax
\newcommand{\GL}{\operatorname{GL}}
\let\St\relax
\newcommand{\St}{\operatorname{St}}
\let\GV\relax
\newcommand{\GV}{\operatorname{GV}}
\let\E\relax
\newcommand{\E}{\operatorname{E}}
\let\Diff\relax
\newcommand{\Diff}{\operatorname{Diff}}
\let\Symp\relax
\newcommand{\Symp}{\operatorname{Symp}}
\let\Ham\relax
\newcommand{\Ham}{\operatorname{Ham}}
\let\Spec\relax
\newcommand{\Spec}{\operatorname{Spec}}
\let\Ker\relax
\newcommand{\Ker}{\operatorname{Ker}}
\let\vol\relax
\newcommand{\vol}{\operatorname{vol}}
\let\diam\relax
\newcommand{\diam}{\operatorname{diam}}
\let\area\relax
\newcommand{\area}{\operatorname{area}}
\let\Image\relax
\newcommand{\Image}{\operatorname{Im}}
\let\Lip\relax
\newcommand{\Lip}{\operatorname{Lip}}
\let\id\relax
\newcommand{\id}{\operatorname{id}}
\let\supp\relax
\newcommand{\supp}{\operatorname{supp}}
\let\const\relax
\newcommand{\const}{\operatorname{const}}

\chapter{J\"{u}rgen Moser (1994/95)}

\begin{quote}
\it ``Moser's work is distinguished by a rare combination of deep analytical power, geometric insight, and profound contributions to dynamical systems. His discoveries in KAM theory, the implicit function theorem that bears his name with Nash, and his fundamental contributions to the theory of partial differential equations, symplectic geometry, and celestial mechanics have reshaped large areas of mathematics. His proofs are models of elegance and technical mastery.'' --- Wolf Prize Citation
\end{quote}

J\"{u}rgen Kurt Moser (July 4, 1928 -- December 17, 1999) was one of the most versatile and influential mathematicians of the 20th century. The 1994/95 Wolf Prize in Mathematics (shared with Pierre Deligne) recognized his monumental contributions to dynamical systems, the theory of partial differential equations, symplectic geometry, and celestial mechanics. Moser's name is attached to a remarkable array of concepts and theorems across pure and applied mathematics: the KAM (Kolmogorov--Arnold--Moser) theorem, the Nash--Moser implicit function theorem, the Moser twist theorem, the Moser iteration technique, the Moser--Trudinger inequality, the Moser flow on symplectic structures, and the Moser invariant for area-preserving maps.

Moser's mathematical style is characterized by a rare fusion of analytic rigor with geometric intuition. He had a gift for extracting the essential difficulty of a problem---whether it was the small divisor obstructions in celestial mechanics, the loss of derivatives in nonlinear PDEs, or the regularity of weak solutions to elliptic equations---and developing precisely the right technique to overcome it. His work bridges dynamical systems, PDE theory, geometry, and analysis in ways that continue to define research frontiers.

\section{Early Life and Career}

J\"{u}rgen Moser was born on July 4, 1928, in K\"{o}nigsberg, East Prussia (now Kaliningrad, Russia). His father was a neurologist, and the family valued education and intellectual pursuits. The upheavals of World War II forced the Moser family to flee East Prussia, and after the war, Moser settled in G\"{o}ttingen, which was then being rebuilt as a center of mathematical research.

Moser studied at the University of G\"{o}ttingen, the historic home of Gauss, Riemann, Hilbert, and Klein. He completed his PhD in 1952 under the supervision of Franz Rellich, a leading figure in spectral theory and mathematical physics. Moser's doctoral thesis on spectral theory for differential operators already displayed his characteristic blend of analysis and concrete computation.

After his PhD, Moser spent several years at the University of G\"{o}ttingen as a research assistant. In 1955, he moved to the United States, holding positions at the Massachusetts Institute of Technology and, crucially, at the Courant Institute of Mathematical Sciences at New York University. The Courant Institute, under the leadership of Richard Courant and later Peter Lax, was a vibrant center for applied mathematics and partial differential equations. Moser thrived in this environment, developing his deep interests in dynamical systems, celestial mechanics, and nonlinear PDEs.

In 1960, Moser joined the faculty at the Courant Institute as a full professor. He remained there until 1970, when he moved to the Massachusetts Institute of Technology. In 1977, he was appointed a professor at the ETH Z\"{u}rich, where he built a world-renowned research group in analysis and geometry. He returned to the United States in 1980 as a professor at the Courant Institute, where he remained for the rest of his career. He served as the director of the Courant Institute from 1984 to 1988. Moser was awarded numerous honors: he was elected to the National Academy of Sciences in 1973, received the Wolf Prize in 1994/95, and was awarded the Cantor Medal of the German Mathematical Society in 1996. He died on December 17, 1999, in Schwerzenbach, Switzerland.

\section{The Kolmogorov--Arnold--Moser (KAM) Theory}

The KAM theory is one of the crowning achievements of 20th-century mathematics. It addresses a fundamental question in dynamical systems: what happens to quasi-periodic motion (motion on invariant tori) when an integrable Hamiltonian system is slightly perturbed? The question is central to celestial mechanics---the solar system is approximately integrable (planets moving around the sun with weak mutual perturbations), and one wants to know whether its quasi-periodic structure persists over long times.

\subsection{Background: Integrable Systems and Small Divisors}

A Hamiltonian system with $n$ degrees of freedom is \textit{completely integrable} (in the sense of Liouville--Arnold) if there exist $n$ independent integrals of motion in involution. By the Arnold--Liouville theorem, such a system admits \textit{action-angle variables} $(I, \theta) \in D \times \mathbb{T}^n$, where $D \subset \R^n$ is a domain and $\mathbb{T}^n$ is the $n$-torus, in which the Hamiltonian depends only on the actions:
\[
H_0 = H_0(I).
\]
The equations of motion become
\[
\dot{I} = 0, \qquad \dot{\theta} = \omega(I) = \frac{\partial H_0}{\partial I},
\]
so that the actions $I$ are constant and the angles $\theta$ evolve linearly. The system is foliated by invariant $n$-tori $\{I = \const\}$, each carrying quasi-periodic motion with frequency vector $\omega(I)$.

When a small perturbation is added,
\[
H(I,\theta) = H_0(I) + \varepsilon H_1(I,\theta),
\]
the invariant tori are generally destroyed. The fundamental question is whether some of them persist, deformed but present, for sufficiently small $\varepsilon$.

The difficulty lies in the \textit{small divisor problem}. To find invariant tori, one constructs a near-identity canonical transformation $(I,\theta) \mapsto (I',\theta')$ that removes the perturbation to leading order. This involves solving equations of the form
\[
\langle \omega, \partial_\theta \Phi \rangle = \text{known function},
\]
which, when expanded in Fourier series, leads to denominators $\langle k, \omega \rangle$ for $k \in \Z^n \setminus \{0\}$. These denominators can be arbitrarily small---this is the small divisor problem that had frustrated astronomers and mathematicians since Lagrange and Laplace.

\subsection{The KAM Theorem}

The breakthrough came in 1954, when Andrey Kolmogorov announced a theorem at the International Congress of Mathematicians in Amsterdam: under a non-degeneracy condition, most invariant tori persist under small perturbations. Kolmogorov sketched a proof using a rapidly convergent Newton iteration (the ``super-convergent'' method), but the detailed analytic estimates were worked out by Vladimir Arnold (for analytic Hamiltonians) and J\"{u}rgen Moser (for finitely differentiable Hamiltonians).

\begin{theorem}[KAM Theorem, Kolmogorov--Arnold--Moser]
Let $H_0(I)$ be a real-analytic Hamiltonian on a domain $D \subset \R^n$ satisfying the \textit{Kolmogorov non-degeneracy condition}
\[
\det\left(\frac{\partial^2 H_0}{\partial I^2}\right) \neq 0 \qquad \text{on } D.
\]
Let $H(I,\theta) = H_0(I) + \varepsilon H_1(I,\theta)$ be a real-analytic perturbation. Then for sufficiently small $\varepsilon > 0$, there exists a Cantor-like set $\mathcal{C}_\varepsilon \subset D$ of positive measure such that for every $I_0 \in \mathcal{C}_\varepsilon$, the unperturbed torus $\{I = I_0\}$ survives as a deformed invariant torus for the perturbed system. These tori carry quasi-periodic motion with frequencies $\omega(I_0)$ satisfying the \textit{Diophantine condition}
\[
|\langle k, \omega(I_0) \rangle| \geq \frac{\gamma}{|k|^{\tau}} \qquad \text{for all } k \in \Z^n \setminus \{0\},
\]
with $\gamma > 0$ and $\tau > n-1$. The union of the surviving tori has Lebesgue measure approaching $|\mathcal{C}_\varepsilon| \to |D|$ as $\varepsilon \to 0$.
\end{theorem}

The non-degeneracy condition $\det(\partial^2 H_0 / \partial I^2) \neq 0$ ensures that the frequency map $\omega(I) = \partial H_0 / \partial I$ is a local diffeomorphism, so that most tori satisfy the required Diophantine condition. The perturbation $H_1$ can destroy the rational tori (those with resonant frequencies), but the irrational tori, which form a set of full measure, survive.

\subsection{Moser's Contribution: From Analytic to Finite Differentiability}

Arnold's proof of the KAM theorem in 1963 required analyticity of the Hamiltonian. But the three-body problem of celestial mechanics involves functions that are real-analytic only away from collisions, and Moser recognized that a version of the theorem for finitely differentiable systems was essential for physical applications.

Moser's 1962 paper \textit{On invariant curves of area-preserving mappings of an annulus} established the KAM theorem for mappings of class $C^{333}$ (later improved to $C^{3+\delta}$ and eventually to $C^{1+\alpha}$). The key innovation was a smoothing technique: instead of working directly with a finitely differentiable map, Moser smoothed it at each step of the Newton iteration, using convolution with a smooth kernel. The loss of differentiability from the smoothing was compensated by the super-exponential convergence of the Newton method. This idea of \textit{smoothing in the iteration} later became the core of the Nash--Moser implicit function theorem.

Moser's work completed the KAM theory in its full generality. The theorem now stands as one of the principal results of Hamiltonian dynamics. It explains why the solar system is stable over astronomical timescales (the planets' orbits are quasi-periodic with high probability), it accounts for the existence of the asteroid belt's Kirkwood gaps (resonances destroy tori at rational frequencies), and it has found applications to plasma physics, accelerator physics, and molecular dynamics.

\section{The Nash--Moser Implicit Function Theorem}

The classical implicit function theorem requires the derivative of a map to be an isomorphism between Banach spaces at the point of interest. In many problems in differential geometry and PDE theory, however, the natural map between function spaces loses derivatives: the derivative is bounded from a space of higher differentiability to one of lower differentiability, and one cannot directly apply the Banach space inverse function theorem. This is the ``loss of derivatives'' problem.

\subsection{Moser's Insight: Taming the Loss of Derivatives}

In 1956, John Nash stunned the mathematical world by proving the existence of isometric embeddings of Riemannian manifolds into Euclidean space. Nash's proof involved a remarkable iteration scheme that compensated for the loss of derivatives by using smoothing operators. The proof was extraordinarily intricate, and many mathematicians found it difficult to understand.

In 1966, Moser published \textit{A rapidly convergent iteration method and nonlinear partial differential equations}, in which he extracted the essential analytic core of Nash's argument and presented it as a general implicit function theorem. Moser showed that the Newton iteration could be modified to work in Fr\'echet spaces (rather than Banach spaces) by inserting smoothing operators at appropriate steps. The key estimates are ``tame'' estimates that only lose a controlled number of derivatives.

\subsection{Statement of the Theorem}

The Nash--Moser theorem deals with a \textit{graded Fr\'echet space}---a Fr\'echet space $X$ whose topology is defined by an increasing family of norms $\{\|\cdot\|_n\}_{n \in \N}$. A typical example is $C^\infty(M)$, the space of smooth functions on a compact manifold, with the $C^k$ norms. A map $F: X \to Y$ is \textit{tame} if it satisfies estimates of the form
\[
\|F(u)\|_n \leq C(1 + \|u\|_{n+r})
\]
for some constant $r$ (the \textit{loss of derivatives}) and all $n$. The theorem requires the existence of \textit{smoothing operators} $S_t: X \to X$ ($t \geq 1$) that commute with differentiation and satisfy
\[
\|S_t u\|_n \leq C t^{n-m} \|u\|_m \quad (n \geq m), \qquad
\|(I - S_t)u\|_m \leq C t^{-(n-m)} \|u\|_n \quad (m \leq n).
\]

\begin{theorem}[Nash--Moser Implicit Function Theorem]
Let $X$ and $Y$ be graded Fr\'echet spaces admitting smoothing operators, and let $F: U \subset X \to Y$ be a smooth tame map on an open neighborhood $U$ of $0$. Suppose that for each $u \in U$, the Fr\'echet derivative $DF(u): X \to Y$ admits a \textit{tame right inverse}: there exists a family of linear maps $L(u): Y \to X$ such that $DF(u) \circ L(u) = \id_Y$, with tame estimates
\[
\|L(u)v\|_n \leq C(1 + \|u\|_{n+r})\|v\|_{n+r}
\]
for some fixed $r$ independent of $u$. Then $F$ is locally invertible: there exists a neighborhood $V \subset Y$ of $F(0)$ and a tame map $G: V \to X$ such that $F(G(v)) = v$ for all $v \in V$.
\end{theorem}

The theorem is a powerful tool in nonlinear PDE theory. It has been used to:
\begin{itemize}
\item Prove the existence of isometric embeddings of Riemannian manifolds (Nash's original theorem).
\item Construct solutions to the Monge--Amp\`ere equation in K\"ahler geometry (the Calabi--Yau theorem, proved by Yau using different but related methods).
\item Show the existence of quasi-periodic solutions to nonlinear wave and Schr\"odinger equations (the KAM theory for PDEs, developed by Kuksin, Wayne, and others).
\item Establish the existence of minimal surfaces and harmonic maps under various geometric constraints.
\end{itemize}

The Nash--Moser theorem is essential because it bridges the gap between the Banach space implicit function theorem and the reality of PDE problems, where the natural function spaces are Fr\'echet spaces of smooth functions and every differentiation involves a loss of derivatives.

\section{Moser's Twist Theorem}

The twist theorem is a particularly elegant application of KAM theory to two-dimensional area-preserving maps. It concerns maps of the annulus $A = \mathbb{T}^1 \times [-1, 1]$ that are exact area-preserving and satisfy a twist condition: the rotation number (the average angular displacement) varies strictly monotonically with the radial coordinate.

Consider a map $\varphi: A \to A$ of the form
\[
\varphi(\theta, r) = (\theta + \alpha(r), r + \beta(\theta, r)),
\]
where $(\theta, r)$ are coordinates on the annulus, $\alpha(r)$ is the rotation number, and $\varphi$ preserves the area form $d\theta \wedge dr$. The \textit{twist condition} is $\alpha'(r) \neq 0$, meaning that points at different radii rotate at different speeds.

\begin{theorem}[Moser's Twist Theorem]
Let $\varphi$ be a $C^k$ area-preserving twist map of the annulus with $k \geq 5$ (or $k \geq 3$ with additional conditions). If the rotation number $\omega \in \R \setminus \Z$ satisfies a Diophantine condition
\[
|q\omega - p| \geq C|q|^{-\mu} \quad \text{for all } p, q \in \Z,\; q \neq 0,
\]
with $\mu > 1$, then there exists an invariant circle $\Gamma$ of rotation number $\omega$ that is a $C^1$ closed curve transverse to the radial direction. Moreover, $\Gamma$ is the graph of a $C^1$ function $r = u(\theta)$.
\end{theorem}

The twist theorem is of fundamental importance in dynamical systems because it provides a rigorous mechanism for the persistence of invariant circles in near-integrable area-preserving maps. It applies directly to the \textit{standard map}
\[
\varphi(\theta, r) = (\theta + r',\; r' = r + \varepsilon \sin\theta),
\]
a paradigmatic model for the onset of chaos in Hamiltonian dynamics. For small $\varepsilon$, the standard map has many invariant circles (KAM tori); as $\varepsilon$ increases beyond a critical threshold (the Greene criterion), all invariant circles are destroyed and chaotic transport becomes possible.

The twist theorem also provides a rigorous foundation for the \textit{Poincar\'e--Birkhoff fixed point theorem:} any area-preserving twist map of the annulus has at least two fixed points, corresponding to periodic orbits. More generally, the Aubry--Mather theory extends these ideas to study the breakdown of invariant circles and the emergence of Cantor-like invariant sets (Aubry--Mather sets).

\section{The Calabi Invariant and Area-Preserving Maps}

Area-preserving diffeomorphisms (symplectomorphisms in two dimensions) have a rich geometric structure. For the closed unit disk $D^2$, the group $\Diff_{\area}(D^2)$ of area-preserving diffeomorphisms that fix the boundary pointwise has an important numerical invariant, discovered by Eugenio Calabi.

Let $\varphi \in \Diff_{\area}(D^2)$ be an area-preserving diffeomorphism. Since $\varphi$ preserves the area form $dx \wedge dy$, we can write $\varphi$ as the time-1 map of a time-dependent Hamiltonian flow. Equivalently, there exists a 1-form $\alpha$ on $D^2$ such that $\varphi^*(dx \wedge dy) = dx \wedge dy$ and $d\alpha = \varphi^*(dx \wedge dy) - dx \wedge dy = 0$, so $\alpha$ is closed. Because $D^2$ is simply connected, $\alpha$ is exact: $\alpha = dS$ for some function $S$ (the generating function). The \textit{Calabi invariant} is defined as
\[
C(\varphi) = \int_{D^2} S \; dx \wedge dy.
\]

The Calabi invariant is independent of the choice of primitive $S$ up to an additive constant, and it satisfies:
\begin{itemize}
\item $C(\varphi \circ \psi) = C(\varphi) + C(\psi)$ (it is a homomorphism from $\Diff_{\area}(D^2)$ to $\R$).
\item $C(\varphi) = 0$ if $\varphi$ is isotopic to the identity through area-preserving diffeomorphisms with compact support.
\item For a rotation $R_\theta$ of the disk, $C(R_\theta) = \theta \cdot \area(D^2)/2\pi$.
\end{itemize}

Moser made fundamental contributions to the study of the group of area-preserving diffeomorphisms. In a seminal 1965 paper, he proved the following classification theorem:

\begin{theorem}[Moser, 1965]
The group $\Diff_{\area}(D^2)_0$ of area-preserving diffeomorphisms of the disk that are isotopic to the identity is a simple group: it contains no nontrivial normal subgroups. More precisely, the Calabi invariant is the unique nontrivial homomorphism from $\Diff_{\area}(D^2)_0$ to $\R$, and its kernel is a simple group.
\end{theorem}

This result is part of a broader investigation into the structure of diffeomorphism groups that Moser initiated. He showed that the Calabi invariant classifies area-preserving diffeomorphisms up to a certain equivalence, and he established deep connections between the group $\Diff_{\area}(M)$ and the symplectic topology of $M$. These ideas anticipated later developments in the theory of $C^0$ symplectic topology, the flux homomorphism, and the geometry of diffeomorphism groups (the Hofer metric, in particular).

\section{Stability of the Solar System}

The problem of the stability of the solar system has engaged astronomers and mathematicians since Newton. The planets interact gravitationally, and their mutual perturbations cause slow changes in their orbital elements. For two centuries, the question remained unresolved: are the planetary orbits stable over astronomical timescales, or will the solar system eventually dissolve due to cumulative perturbations?

Moser's contribution to this problem was both specific and general. On the specific side, he studied the stability of the Lagrangian triangular equilibrium points in the restricted three-body problem.

\subsection{The Restricted Three-Body Problem}

The restricted three-body problem considers the motion of a massless particle (e.g., an asteroid) under the gravitational influence of two massive bodies (e.g., the Sun and Jupiter) moving in circular orbits about their common center of mass. There are five equilibrium points (Lagrange points): three collinear points ($L_1, L_2, L_3$) and two equilateral triangle points ($L_4, L_5$). The triangular points form equilateral triangles with the two primaries.

The stability of $L_4$ and $L_5$ depends on the mass ratio of the two primaries. For the planar restricted problem, linear stability (stable in the linearized approximation) holds when the mass ratio satisfies $\mu / (1-\mu) < 1/27$, i.e., when the smaller primary is less than about $3.7\%$ of the total mass. For the Sun--Jupiter system, the mass ratio is about $0.1\%$, well within the stable range.

\begin{theorem}[Moser, 1958]
For the planar restricted three-body problem with the mass ratio satisfying the stability condition $\mu < \mu_0$ (the Routh criterion), the triangular Lagrange points $L_4$ and $L_5$ are stable in the sense of Lyapunov for almost all initial conditions sufficiently close to equilibrium. More precisely, the set of initial conditions that do not lead to stable quasi-periodic motion has measure zero.
\end{theorem}

Moser's proof used KAM theory: near $L_4$ and $L_5$, the Hamiltonian expands as $H = H_2 + H_3 + \cdots$, where $H_2$ is a harmonic oscillator. The frequencies of $H_2$ satisfy the non-resonance condition under the Routh stability criterion. By applying the KAM theorem, Moser showed that most trajectories near the equilibrium lie on invariant tori and are therefore confined for all time.

This result has a concrete astronomical application: it explains the existence of the \textit{Trojan asteroids}, which orbit at the $L_4$ and $L_5$ points of the Sun--Jupiter system. Moser's theorem guarantees that these asteroids, even when perturbed by other planets, remain near the Lagrange points over astronomical timescales.

On the general side, Moser's work on the stability of the solar system contributed to the broader KAM framework that explains why nearly integrable Hamiltonian systems exhibit persistent quasi-periodic motion. While the full $n$-body problem remains incompletely resolved (the question of the long-term stability of the inner planets is still an active area of research), Moser's results provide a rigorous foundation for believing that the solar system is stable over extremely long---though perhaps not infinite---timescales.

\section{Integrable Hamiltonian Systems}

Moser made several deep contributions to the theory of integrable systems. While the KAM theorem concerns the persistence of integrability under perturbation, Moser also studied the structure of exactly integrable systems and the mechanisms by which integrability arises.

\subsection{The Liouville--Arnold Theorem}

A Hamiltonian system with $n$ degrees of freedom is completely integrable if there exist $n$ independent functions $F_1, \dots, F_n$ in involution:
\[
\{F_i, F_j\} = 0 \quad \text{for all } i,j.
\]
The Liouville--Arnold theorem states that the level sets of the $F_i$ are (when compact and connected) $n$-tori, and the motion is conjugate to linear flow on these tori. Moser's work clarified the regularity of the action-angle variables and established sharp conditions under which the conjugation is smooth.

\subsection{Integrable Systems and Spectral Theory}

Moser made important contributions to the theory of finite-dimensional integrable systems arising from the spectral theory of linear operators (the Lax pair formalism). In particular, he studied integrable systems associated with the Toda lattice and the periodic Korteweg--de Vries (KdV) equation. The periodic KdV equation can be viewed as an infinite-dimensional integrable Hamiltonian system, and its finite-gap solutions correspond to integrable systems on the Jacobian varieties of Riemann surfaces.

Moser's 1975 paper \textit{Finite-dimensional integrable Hamiltonian systems related to isospectral deformations of Jacobi matrices} established a deep connection between: (i) the theory of orthogonal polynomials, (ii) classical continued fractions, (iii) the motion of particles under an exponential interaction (the Toda lattice), and (iv) the spectral theory of Jacobi matrices. This paper is a masterpiece of integration between analysis, algebra, and geometry.

\subsection{The Toda Lattice}

The Toda lattice is a system of particles on a line with nearest-neighbor exponential interactions:
\[
H(p,q) = \frac{1}{2}\sum_{i=1}^N p_i^2 + \sum_{i=1}^{N-1} e^{q_i - q_{i+1}}.
\]

Moser showed that the Toda lattice is completely integrable by constructing $N$ independent integrals of motion in involution. The key observation is that the equations of motion can be written in Lax form $\dot{L} = [B, L]$, where $L$ is a Jacobi matrix. The eigenvalues of $L$ are constants of motion, and the isospectral flow of $L$ is the Toda flow. Moser's analysis connected the Toda lattice to the theory of orthogonal polynomials, the moment problem, and the spectral theory of second-order difference operators.

This work opened up a vast area of research at the intersection of integrable systems, spectral theory, and algebraic geometry. The Toda lattice became a prototype for finite-dimensional integrable systems derived from infinite-dimensional PDEs (the KdV hierarchy, in particular), and Moser's spectral approach provided the template for understanding the algebro-geometric solutions of these equations.

\section{The Moser Iteration Technique}

One of Moser's most widely used contributions is the iteration technique he developed for proving regularity of weak solutions to elliptic partial differential equations. The technique appears in Moser's 1960--1961 papers on the Harnack inequality and is a cornerstone of elliptic PDE theory.

\subsection{The Problem of Regularity}

Consider a second-order elliptic equation in divergence form:
\[
-\sum_{i,j=1}^n \partial_i(a_{ij}(x)\partial_j u) = 0,
\]
where the coefficients $a_{ij}(x)$ are measurable and satisfy the ellipticity condition
\[
\lambda^{-1}|\xi|^2 \leq \sum_{i,j} a_{ij}(x) \xi_i \xi_j \leq \lambda |\xi|^2
\]
for some $\lambda > 0$. A weak solution $u \in H^1(\Omega)$ satisfies
\[
\int_\Omega \sum_{i,j} a_{ij} \partial_i u \, \partial_j \varphi \, dx = 0
\]
for all test functions $\varphi \in C_c^\infty(\Omega)$. The fundamental question is: how regular are weak solutions?

De Giorgi (1957) proved that weak solutions are H\"older continuous for $n \geq 3$, and Nash (1958) proved the Harnack inequality for parabolic equations. Moser, in a series of papers (1960, 1961, 1964), gave a new, elegant proof of the H\"older continuity and the Harnack inequality for elliptic equations with measurable coefficients.

\subsection{Moser's Iteration}

The Moser iteration technique proceeds as follows. Let $u \geq 0$ be a weak solution. For any $p > 0$ and any cutoff function $\zeta$, one applies the weak formulation to test functions involving powers $u^{p-1} \zeta^2$. Through careful estimation using the Sobolev and H\"older inequalities, one derives an iterative inequality of the form
\[
\|u\|_{L^{p_{k+1}}(B_{r_{k+1}})} \leq C^{1/p_k} \|u\|_{L^{p_k}(B_{r_k})},
\]
where $p_{k+1} = \kappa p_k$ with $\kappa = n/(n-2) > 1$, and the balls $B_{r_k}$ shrink to a point. Iterating this inequality, one obtains an $L^\infty$ bound on $u$:
\[
\sup_{B_{R/2}} |u| \leq C \|u\|_{L^2(B_R)}.
\]

Combining this with the Harnack inequality (which Moser also proved via iteration), one obtains:
\begin{itemize}
\item Every weak solution is H\"older continuous.
\item Non-negative weak solutions satisfy the Harnack inequality:
\[
\sup_{B_{R/2}} u \leq C \inf_{B_{R/2}} u,
\]
with the constant $C$ depending only on $n$, $\lambda$, and $R$.
\end{itemize}

\begin{theorem}[Moser's Harnack Inequality]
Let $u \geq 0$ be a weak solution of an elliptic equation in divergence form with bounded measurable coefficients satisfying the uniform ellipticity condition. Then for any ball $B_R \subset \Omega$,
\[
\sup_{B_{R/2}} u \leq C \inf_{B_{R/2}} u,
\]
where $C = C(n, \lambda)$ is independent of $u$.
\end{theorem}

The Moser iteration technique has become a standard tool in PDE theory. It is used in:
\begin{itemize}
\item The regularity theory of elliptic and parabolic equations.
\item The proof of $L^p$ estimates for solutions of the heat equation.
\item The study of nonlinear elliptic equations (the $p$-Laplacian, minimal surface equation).
\item The analysis of the Yamabe equation and geometric PDEs on manifolds.
\item The theory of harmonic maps between Riemannian manifolds.
\end{itemize}

\section{The Moser--Trudinger Inequality}

The Sobolev embedding theorem states that for $p < n$, the space $W^{1,p}_0(\Omega)$ embeds continuously into $L^{p^*}(\Omega)$ where $p^* = np/(n-p)$. For the borderline case $p = n$, the critical exponent is $p^* = \infty$, but $W^{1,n}_0(\Omega)$ does not embed into $L^\infty(\Omega)$. The question is: what is the optimal exponential integrability of functions in $W^{1,n}_0(\Omega)$?

Trudinger (1967) proved that $W^{1,n}_0(\Omega)$ embeds into the Orlicz space associated with the exponential function: for any $u \in W^{1,n}_0(\Omega)$, the integral $\int_\Omega \exp(\alpha |u|^{n/(n-1)})\,dx$ is finite for any $\alpha > 0$. However, this did not identify the optimal (largest) $\alpha$ for which a uniform bound holds.

Moser (1971) solved this problem by determining the sharp constant.

\begin{theorem}[Moser--Trudinger Inequality]
Let $\Omega \subset \R^n$ be a bounded domain and let
\[
\alpha_n = n \omega_{n-1}^{1/(n-1)},
\]
where $\omega_{n-1}$ is the surface measure of the unit sphere $S^{n-1} \subset \R^n$. For any $u \in W^{1,n}_0(\Omega)$ with $\int_\Omega |\nabla u|^n\,dx \leq 1$, we have
\[
\int_\Omega \exp\bigl(\alpha_n |u|^{n/(n-1)}\bigr)\,dx \leq C_n |\Omega|,
\]
where $C_n$ is a constant depending only on $n$. The constant $\alpha_n$ is sharp: for any $\alpha > \alpha_n$, there exists a sequence of functions $u_k \in W^{1,n}_0(\Omega)$ with $\int_\Omega |\nabla u_k|^n\,dx \leq 1$ such that
\[
\int_\Omega \exp\bigl(\alpha |u_k|^{n/(n-1)}\bigr)\,dx \to \infty \quad \text{as } k \to \infty.
\]

In particular, for $n = 2$, we have $\alpha_2 = 4\pi$, and the inequality reads
\[
\int_\Omega \exp(4\pi u^2)\,dx \leq C |\Omega| \quad \text{whenever} \quad \int_\Omega |\nabla u|^2\,dx \leq 1,
\]
with $4\pi$ the optimal constant.
\end{theorem}

The Moser--Trudinger inequality is of fundamental importance in geometric analysis. It appears in:
\begin{itemize}
\item The theory of surfaces of constant mean curvature (the study of the Willmore functional).
\item The analysis of the prescribed Gaussian curvature equation on $S^2$.
\item The study of the Chern--Simons gauge theory in two dimensions.
\item The bubbling analysis for harmonic maps from surfaces.
\item The function theory on Riemann surfaces (the Onofri inequality, which is a sharp form of the Moser--Trudinger inequality on the sphere).
\end{itemize}

Moser's proof used symmetrization techniques and careful analysis of the exponential integral for radial functions. The sharpness relies on constructing a sequence of functions that concentrate near a point---a precursor to the bubbling phenomenon in geometric analysis.

\section{The Moser Flow on Symplectic Structures}

In symplectic geometry, the question of when two symplectic forms on a manifold are equivalent arises naturally. Moser's 1965 paper \textit{On the volume of a symplectic manifold} introduced a powerful technique---now called the \textit{Moser trick} or \textit{Moser flow}---that has become one of the most fundamental tools in symplectic geometry.

\subsection{The Moser Stability Theorem}

Let $M$ be a compact manifold without boundary. Let $\omega_0$ and $\omega_1$ be two symplectic forms on $M$ that are cohomologous: $[\omega_0] = [\omega_1] \in H^2(M; \R)$. Consider the linear interpolation $\omega_t = (1-t)\omega_0 + t\omega_1$. Since the forms are cohomologous, we have $\omega_1 - \omega_0 = d\alpha$ for some 1-form $\alpha$.

Moser's idea is to look for a time-dependent vector field $X_t$ whose flow $\varphi_t$ satisfies $\varphi_t^*\omega_t = \omega_0$. Differentiating this equation gives
\[
\varphi_t^*\bigl(L_{X_t}\omega_t + \dot{\omega}_t\bigr) = 0,
\]
where $L_{X_t}$ is the Lie derivative. By the Cartan formula, $L_{X_t}\omega_t = d(\iota_{X_t}\omega_t)$ (since $d\omega_t = 0$). Thus we need to solve the equation
\[
d(\iota_{X_t}\omega_t) + \dot{\omega}_t = 0,
\]
which reduces to
\[
\iota_{X_t}\omega_t = -\alpha,
\]
since $\dot{\omega}_t = \omega_1 - \omega_0 = d\alpha$. Because $\omega_t$ is non-degenerate for all $t$ (the space of symplectic forms is convex), this equation uniquely determines $X_t$ from $\alpha$. The flow $\varphi_t$ of $X_t$ then satisfies $\varphi_t^*\omega_t = \omega_0$, and in particular $\varphi_1^*\omega_1 = \omega_0$.

\begin{theorem}[Moser's Stability Theorem]
Let $M$ be a compact manifold without boundary. If $\omega_0$ and $\omega_1$ are cohomologous symplectic forms on $M$, then there exists a diffeomorphism $\varphi$ isotopic to the identity such that $\varphi^*\omega_1 = \omega_0$.
\end{theorem}

\subsection{Consequences and Generalizations}

The Moser flow has several important consequences:

\begin{itemize}
\item \textbf{Darboux's theorem:} Every symplectic form is locally diffeomorphic to the standard symplectic form on $\R^{2n}$. This follows by applying the Moser trick to a family of forms that interpolate between $\omega$ and the standard form in a Darboux chart.

\item \textbf{Classification of symplectic forms on surfaces:} On a closed orientable surface $\Sigma_g$, two symplectic forms are equivalent if and only if they have the same total volume.

\item \textbf{The Weinstein conjectures:} Moser's ideas influenced Weinstein's program on the geometry of symplectic manifolds and the Arnold conjecture on the number of fixed points of Hamiltonian diffeomorphisms.

\item \textbf{The Moser trick for contact forms:} An analogous result holds in contact geometry: cohomologous contact forms on a compact manifold are contactomorphic.

\item \textbf{The Neishtadt--Moser theorem on adiabatic invariants:} In Hamiltonian systems with slowly varying parameters, the action variables are approximate integrals of motion. Moser's work on adiabatic invariants used similar deformation techniques.
\end{itemize}

\subsection{The Volume Theorem}

A related result, also proved by Moser in the same paper, concerns the topology of symplectic forms:

\begin{theorem}[Moser, 1965]
On a compact orientable manifold $M^{2n}$, two volume forms $\mu_0$ and $\mu_1$ with the same total volume $\int_M \mu_0 = \int_M \mu_1$ are diffeomorphic: there exists a diffeomorphism $\varphi$ isotopic to the identity such that $\varphi^*\mu_1 = \mu_0$.
\end{theorem}

This theorem shows that volume forms, unlike symplectic forms, have no local invariants beyond the total volume. The proof uses the same Moser flow technique: given a family $\mu_t = (1-t)\mu_0 + t\mu_1$, one constructs a vector field $X_t$ whose flow pushes $\mu_0$ to $\mu_t$.

\section{Impact and Legacy}

J\"{u}rgen Moser's work has had a transformative impact on mathematics, and his ideas continue to shape research across multiple fields.

\subsection{Dynamical Systems and Celestial Mechanics}

The KAM theory, of which Moser was a principal architect, is one of the most important results in dynamical systems. It has been extended and generalized in countless directions: to infinite-dimensional systems (KAM for PDEs, developed by Kuksin, Wayne, and Eliasson), to lower-dimensional tori, to quasi-periodic solutions in Hamiltonian systems, and to the study of Arnold diffusion. The KAM theory provides the foundation for understanding stability in nearly integrable systems and has applications to particle accelerators, plasma confinement in fusion devices, and the orbital dynamics of exoplanets.

\subsection{Partial Differential Equations}

The Moser iteration technique and the Moser--Trudinger inequality are standard tools in PDE analysis. The iteration technique is used wherever one needs to upgrade integrability to pointwise bounds---a common step in regularity theory. The Moser--Trudinger inequality is essential in the analysis of critical growth problems, particularly in two-dimensional conformal geometry.

\subsection{Symplectic Geometry}

The Moser flow and Moser stability theorem are among the first results that a student of symplectic geometry encounters. They underpin the foundational results of the field (Darboux's theorem, the Weinstein tubular neighborhood theorem, and the classification of symplectic forms). The Moser trick is a paradigm for constructing diffeomorphisms with desired properties by solving linear equations.

\subsection{Integrable Systems and Spectral Theory}

Moser's work on the Toda lattice and finite-dimensional integrable systems opened up the connection between integrable PDEs, spectral theory, and algebraic geometry. These ideas have been extended to the study of completely integrable systems on Lie algebras (the Kostant--Adler--Symes construction), the analysis of the KdV equation, and the theory of Riemann surfaces.

\subsection{Expository Writing and Education}

Moser was a superb expositor. His lecture notes and survey articles are models of clarity. His book \textit{Stable and Random Motions in Dynamical Systems} (1973, Princeton University Press, Annals of Mathematics Studies 77) is a classic introduction to KAM theory and the foundations of Hamiltonian dynamics. His article \textit{On the volume of a symplectic manifold} is among the most cited papers in symplectic geometry. His collaboration with Charles Siegel on the theory of Riemann surfaces produced influential lecture notes.

Moser supervised many PhD students who became leading mathematicians, including Henry McKean, Richard Laver, and others. He was known for his generosity, his broad mathematical curiosity, and his insistence on clarity and rigor.

\subsection{Continuing Influence}

Moser's legacy endures in the many areas of mathematics that he touched. The KAM theory remains an active area of research, with recent advances focusing on the breakdown of invariant tori, the transition to chaos, and the extension to infinite-dimensional Hamiltonian systems. The Nash--Moser theorem continues to be refined and applied to new problems in geometry and PDEs. The Moser--Trudinger inequality and its generalizations (the Adams inequality, the Onofri inequality, the Beckner inequality) are central to conformal geometry and geometric analysis. The Moser flow remains an essential tool in symplectic and contact geometry.

Moser's unique ability to combine deep analysis with geometric insight, and his talent for identifying the essential mathematical structure underlying physical problems, places him among the great mathematicians of the 20th century.

\section{Timeline}

\begin{tabular}{|l|l|}
\hline
\textbf{Year} & \textbf{Event} \\ \hline
1928 & Born on July 4 in K\"{o}nigsberg, East Prussia \\ \hline
1940s & Postwar studies at the University of G\"{o}ttingen \\ \hline
1952 & Ph.D.\ from the University of G\"{o}ttingen (advisor: Franz Rellich) \\ \hline
1953--1955 & Assistant at G\"{o}ttingen; early work on spectral theory \\ \hline
1955 & Moved to the United States; positions at MIT and later the Courant Institute \\ \hline
1958 & Proof of the stability of the Lagrange triangular equilibrium points \\ \hline
1960--1961 & Moser iteration technique and the Harnack inequality for elliptic PDEs \\ \hline
1962 & KAM theorem for $C^{333}$ area-preserving maps \\ \hline
1965 & Moser stability theorem (the Moser flow) for symplectic forms \\ \hline
1966 & Nash--Moser implicit function theorem published \\ \hline
1970 & Moved to MIT as a full professor \\ \hline
1971 & Sharp Moser--Trudinger inequality (optimal constant in exponential Sobolev embedding) \\ \hline
1975 & Work on finite-dimensional integrable systems and the Toda lattice \\ \hline
1977 & Appointed professor at the ETH Z\"{u}rich \\ \hline
1980 & Returned to the Courant Institute, New York University \\ \hline
1984--1988 & Director of the Courant Institute of Mathematical Sciences \\ \hline
1994/95 & Awarded the Wolf Prize in Mathematics (shared with Pierre Deligne) \\ \hline
1996 & Cantor Medal of the German Mathematical Society \\ \hline
1999 & Died on December 17 in Schwerzenbach, Switzerland \\ \hline
\end{tabular}

\section*{Exercises}
\addcontentsline{toc}{section}{Exercises}

\begin{enumerate}
\item Let $\omega = (\omega_1, \dots, \omega_n)$ be a vector of frequencies. Show that if the $\omega_i$ and $1$ are linearly independent over $\Q$ (i.e., $\langle k, \omega \rangle = 0$ only for $k = 0$), then the linear flow $\theta \mapsto \theta + \omega t$ on $\mathbb{T}^n$ is dense. Show that the Diophantine condition $|\langle k, \omega \rangle| \geq \gamma/|k|^\tau$ implies that the flow is also uniquely ergodic.

\item Prove that the set of numbers $\omega \in [0,1]$ satisfying the Diophantine condition $|q\omega - p| \geq c/|q|^{\mu}$ for all $p/q \in \Q$ has full Lebesgue measure when $\mu > 1$. (This is a key step in the KAM theorem: most frequencies are Diophantine.)

\item Show that the golden ratio $\phi = (1+\sqrt{5})/2$ satisfies the Diophantine condition with $\mu = 1$: $|q\phi - p| \geq 1/(\sqrt{5}\,q)$ for all $p,q \in \Z$,$q \neq 0$. (This is a classical result in Diophantine approximation related to the continued fraction expansion of $\phi$.)

\item Consider the harmonic oscillator $H(p,q) = \frac{1}{2}(p^2 + q^2)$. Find action-angle variables $(I,\theta)$ for this system and verify that the Hamiltonian depends only on $I$. Compute the frequency $\omega(I) = \partial H/\partial I$ and check the Kolmogorov non-degeneracy condition.

\item For the standard map $\varphi(\theta, r) = (\theta + r',\; r' = r + \varepsilon \sin\theta)$, verify that it preserves the area form $d\theta \wedge dr$ and satisfies the twist condition for small $\varepsilon$. For $\varepsilon = 0$, describe the invariant curves. What happens to these curves as $\varepsilon$ increases?

\item Let $F: B_0 \to B_0$ be a $C^1$ map between Banach spaces with $F(0) = 0$ and $DF(0)$ an isomorphism. Prove the standard implicit function theorem using Newton iteration: set $x_{n+1} = x_n - DF(0)^{-1}F(x_n)$ and show that $x_n \to 0$ quadratically. Why does this iteration fail when the derivative loses derivatives?

\item Let $u \in C^\infty(\mathbb{T}^n)$ and define the smoothing operator $S_t$ by truncating the Fourier expansion at frequency $|k| \leq t$: $S_t u(\theta) = \sum_{|k| \leq t} \hat{u}(k) e^{i\langle k,\theta \rangle}$. Show that $\|S_t u\|_{C^m} \leq C t^{m - k} \|u\|_{C^k}$ for $m \geq k$ and $\|(I - S_t) u\|_{C^k} \leq C t^{-(m-k)} \|u\|_{C^m}$ for $k \leq m$. These are the smoothing estimates used in the Nash--Moser theorem.

\item Let $\varphi$ be an area-preserving diffeomorphism of the unit disk $D^2$ that equals a rotation $R_\theta$ near the boundary. Show that the Calabi invariant $C(\varphi) = \theta \cdot \area(D^2)/2\pi$ by constructing the generating function $S$ such that $dS = \varphi^*(dx \wedge dy) - dx \wedge dy$.

\item In the restricted three-body problem, verify that the Lagrange triangular point $L_4$ is an equilibrium when the three bodies form an equilateral triangle. Linearize the equations of motion around $L_4$ and compute the frequencies. Show that these frequencies are real (implying linear stability) when the mass ratio satisfies $\mu/(1-\mu) < 1/27$.

\item Consider the Toda lattice with $N=2$ particles. Write the equations of motion, find a Lax pair $(L,B)$, and verify that $L$ is a $2 \times 2$ symmetric matrix. Show that the eigenvalues of $L$ are constants of motion and that the system is completely integrable.

\item Let $u \in C^2(B_R)$ be a non-negative solution of the Laplace equation $\Delta u = 0$ in $B_R \subset \R^n$. Prove the classical Harnack inequality
\[
\sup_{B_{R/2}} u \leq C \inf_{B_{R/2}} u
\]
using the mean value property. Compare this proof with Moser's iteration technique for elliptic equations with measurable coefficients.

\item Let $\Omega \subset \R^2$ be the unit disk. Consider the family of radial functions $u_\varepsilon(r) = (\log(1/\varepsilon))^{-1/2} \cdot (-\log r)^{1/2}$ for $0 < r < 1$ and $u_\varepsilon \equiv 0$ for $r \leq \varepsilon$. Show that $\int_\Omega |\nabla u_\varepsilon|^2 \, dx = 1$ and compute $\int_\Omega \exp(4\pi u_\varepsilon^2) \, dx$. Examine the limit as $\varepsilon \to 0$ to see why $4\pi$ is the sharp constant in the Moser--Trudinger inequality.

\item Show that the standard symplectic form $\omega_0 = dx_1 \wedge dy_1 + \cdots + dx_n \wedge dy_n$ on $\R^{2n}$ and the form $\omega_1 = 2\omega_0$ are not symplectomorphic on $\R^{2n}$ via a diffeomorphism isotopic to the identity, but are symplectomorphic on $\R^{2n} \setminus \{0\}$. This illustrates the compactness hypothesis in Moser's stability theorem.

\item Let $\omega_0$ and $\omega_1$ be two symplectic forms on $\mathbb{T}^2$ (the 2-torus) such that $\int_{\mathbb{T}^2} \omega_0 = \int_{\mathbb{T}^2} \omega_1 = 1$. Show that $\omega_0$ and $\omega_1$ are symplectomorphic. This is a special case of Moser's theorem: on a closed surface, the only symplectic invariant is the total volume.

\item Prove that the harmonic oscillator $H = \frac{1}{2}(p^2 + \omega^2 q^2)$ has the frequencies $\partial^2 H / \partial I^2 = 0$ when expressed in action-angle variables. That is, the harmonic oscillator is a degenerate integrable system. Explain why the KAM theorem fails for the harmonic oscillator and how the isoenergetic non-degeneracy condition $\det(\partial^2 H / \partial I^2) \neq 0$ is essential.

\end{enumerate}

\section*{Glossary}
\addcontentsline{toc}{section}{Glossary}

\begin{description}
\item[Action-angle variables] A canonical coordinate system $(I,\theta)$ for an integrable Hamiltonian system in which the Hamiltonian depends only on the actions $I$, so that the actions are constant and the angles evolve linearly.

\item[Area-preserving map] A diffeomorphism $\varphi$ of a 2-dimensional manifold such that $\varphi^*\omega = \omega$, where $\omega$ is the area form. Equivalently, $\varphi$ preserves the area of every measurable set.

\item[Calabi invariant] A numerical invariant of an area-preserving diffeomorphism of the disk, defined by integrating the generating function over the disk.

\item[Diophantine condition] A condition on a frequency vector $\omega \in \R^n$ that requires $|\langle k, \omega \rangle| \geq \gamma/|k|^\tau$ for all nonzero integer vectors $k$, ensuring that small denominators are not too small.

\item[Invariant torus] A submanifold of phase space that is invariant under the flow of a Hamiltonian system. In integrable systems, the phase space is foliated by invariant $n$-tori.

\item[KAM theory] The Kolmogorov--Arnold--Moser theory, which states that most invariant tori of an integrable Hamiltonian system persist under small perturbations, provided the frequency satisfies a Diophantine condition.

\item[Kolmogorov non-degeneracy] The condition $\det(\partial^2 H_0/\partial I^2) \neq 0$, which ensures that the frequency map is a local diffeomorphism.

\item[Lagrange points] Five equilibrium points in the restricted three-body problem, of which $L_4$ and $L_5$ form equilateral triangles with the two gravitating bodies and exhibit stability for certain mass ratios.

\item[Lax pair] A pair of matrices $(L,B)$ such that the equation $\dot{L} = [B, L]$ (the Lax equation) is equivalent to an integrable Hamiltonian system. The eigenvalues of $L$ are constants of motion.

\item[Liouville--Arnold theorem] A theorem stating that a completely integrable Hamiltonian system with $n$ degrees of freedom has invariant $n$-tori in phase space, with the motion given by linear flow on these tori.

\item[Moser flow] The flow of a time-dependent vector field constructed to deform one symplectic (or volume) form into another, used to prove Moser's stability theorem.

\item[Moser iteration] A technique for proving regularity of solutions to elliptic PDEs by iterating $L^p$ estimates to obtain $L^\infty$ bounds. The key is the iterative inequality $\|u\|_{L^{p_{k+1}}} \leq C^{1/p_k} \|u\|_{L^{p_k}}$ with $p_k \to \infty$.

\item[Moser--Trudinger inequality] The sharp exponential Sobolev inequality: for $u \in W^{1,n}_0(\Omega)$ with $\|\nabla u\|_{L^n} \leq 1$, the integral $\int_\Omega \exp(\alpha_n |u|^{n/(n-1)})$ is bounded by $C|\Omega|$, and $\alpha_n$ is optimal.

\item[Nash--Moser theorem] An implicit function theorem for smooth tame maps between graded Fr\'echet spaces with smoothing operators, designed to handle problems where derivatives are lost.

\item[Restricted three-body problem] The problem of the motion of a massless particle under the gravitational influence of two massive bodies (primaries) in circular orbits about their common center of mass.

\item[Small divisor problem] The difficulty in perturbation theory arising from denominators $\langle k, \omega \rangle$ that can be arbitrarily small for integer vectors $k$. This problem obstructs the convergence of formal power series in celestial mechanics.

\item[Smoothing operator] An operator $S_t$ that maps a function to a smoother function, typically by truncating its Fourier expansion or convolving with an approximate identity, with controlled estimates in terms of $t$.

\item[Symplectic form] A closed, non-degenerate differential 2-form $\omega$ on an even-dimensional manifold. Darboux's theorem states that every symplectic form is locally equivalent to the standard form $\sum dx_i \wedge dy_i$.

\item[Tame map] A map $F$ between graded Fr\'echet spaces that satisfies $\|F(u)\|_n \leq C(1 + \|u\|_{n+r})$ for some fixed $r$ (the loss of derivatives), with analogous estimates for its derivatives.

\item[Toda lattice] An integrable system of particles on a line with nearest-neighbor exponential interactions. Its equations of motion can be written in Lax form with a Jacobi matrix.

\item[Twist condition] The condition $\alpha'(r) \neq 0$ for an area-preserving map of the annulus, meaning that the rotation number varies strictly with the radial coordinate. This condition ensures the applicability of the twist theorem and the Poincar\'e--Birkhoff theorem.

\item[Weak solution] A solution of a partial differential equation defined via integration by parts, allowing the equation to be satisfied in a distributional sense for measurable coefficients and non-smooth solutions.

\end{description}
